\chapter{Source digests III: monographs and exact examples}
\label{ch:digests-monographs}

This chapter digests two monographs that the approaches to MLC take as
background: Milnor's introductory lectures~\cite{Milnor} and Lyubich's notes
\cite{LyubichNotes}.  It then records exact Julia-set examples that the
monographs use, together with their formal counterparts.  As before, the
statements are those of the sources, up to notation; theorem numbers are
those printed in the sources.

\section{Milnor's introductory lectures}
\label{sec:digest-milnor}

The lectures~\cite{Milnor} are the standard entry point for the quadratic
family.  The following definitions and results concern the quadratic family.

\subsection*{Definitions}

\begin{definition}[Fatou and Julia sets]
\label{def:milnor-fatou-julia}
Let $f$ be a nonconstant holomorphic self-map of the Riemann sphere.  A point
$z_0$ is in the Fatou set if some neighbourhood of $z_0$ carries a normal family of
iterates.  The Julia set $J(f)$ is the complement of the Fatou set; it is closed.
\end{definition}

\begin{definition}[Multipliers and cycle types]
\label{def:milnor-multiplier}
For a cycle $z_1,\dots,z_n$ of period $n$, the multiplier is
$\lambda=(f^{\circ n})'(z_1)=f'(z_1)\cdots f'(z_n)$.  The cycle is attracting,
repelling, or neutral according as $|\lambda|<1$, $|\lambda|>1$, or $|\lambda|=1$,
and superattracting if $\lambda=0$.
\end{definition}

\begin{definition}[Filled Julia set; Böttcher coordinate; external rays]
\label{def:milnor-K-bottcher-rays}
For a polynomial $f$ of degree $d\ge2$ the filled Julia set is
$K(f)=\{z:\ \{f^{\circ n}(z)\}_{n\ge0}\text{ is bounded}\}$.  Near infinity there is a
Böttcher coordinate $\phi$ with $\phi\circ f\circ\phi^{-1}(w)=w^d$, unique up to a
$(d-1)$-st root of unity.  If $K(f)$ is connected, $\phi\colon\C\setminus K\to\C\setminus\overline{\mathbb D}$
is a conformal isomorphism with $\phi(f(z))=\phi(z)^d$.  With $\psi=\phi^{-1}$, the
external ray at angle $t\in\R/\Z$ is $R_t=\psi(\{re^{2\pi it}:r>1\})$, and $f(R_t)=R_{dt}$.
The ray lands at $z_t$ if $\lim_{r\to1}\psi(re^{2\pi it})$ exists.
\end{definition}

\begin{definition}[Quadratic family and parameter sets]
\label{def:milnor-quadratic}
Every quadratic polynomial is affinely conjugate to a unique $f_c(z)=z^2+c$.  Put
\begin{align*}
 \M&=\{c:\ K(f_c)\text{ contains the critical point }0\},\\
 Q&=\{c:\ K(f_c)\text{ has an interior point}\},
\end{align*}
and let $\mathcal H$ be the set of parameters with an attracting periodic orbit in
the finite plane.  Then $\mathcal H\subset Q\subset\M$.  Milnor does not give a
separate definition of ``hyperbolic component''; the phrase refers to components
of the interior of $\M$.
\end{definition}

\subsection*{Results}

\begin{theorem}[Milnor, Theorem~3.3]
\label{thm:milnor-3-3}
Every attracting periodic orbit lies in the Fatou set.  The whole basin of
attraction of such an orbit lies in the Fatou set; its boundary lies in the Julia
set, and every repelling periodic orbit lies in the Julia set.  Status: proved in
the text.
\end{theorem}

\begin{theorem}[Milnor, Lemma~17.1]
\label{thm:milnor-17-1}
$K$ is compact.  It consists of $J$ together with the bounded components of
$\C\setminus J$; these components are simply connected, and $J=\partial K$.
Status: proved in the text.
\end{theorem}

\begin{theorem}[Milnor, Theorem~18.1; cited]
\label{thm:milnor-18-1}
If $K(f)$ is connected, every periodic external ray lands at a periodic point that
is repelling or parabolic.  Status: attributed to Sullivan, Douady and Hubbard; not
proved in the text.
\end{theorem}

\begin{theorem}[Milnor, Theorem~18.2; cited]
\label{thm:milnor-18-2}
If $K(f)$ is connected, every repelling or parabolic periodic point is the landing
point of at least one external ray, which is periodic.  Status: attributed to
Douady and Yoccoz; only the parabolic case is proved in the text.
\end{theorem}

\begin{remark}[What Milnor reports about $\M$]
The text reports, as known from Douady and Hubbard, that $\M$ is connected with
connected complement, that each hyperbolic component of the interior of $\M$ admits a
canonical parametrization, and that the boundary of $\M$ can be studied with external
rays.  It states as widely believed that $\mathcal H=\operatorname{int}\M$ and
$\M=\overline{\mathcal H}$.  It does not state local connectivity of $\M$ as a
conjecture, and no theorem asserting it appears there.
\end{remark}

\section{Lyubich's notes (ICTP, 2013)}
\label{sec:digest-lyubich-notes}

The notes~\cite{LyubichNotes} divide into four parts: conformal and quasiconformal
geometry (including Teichmüller theory and holomorphic motions); the dynamical
plane of the quadratic family (periodic motions, hyperbolic, parabolic and
Misiurewicz maps, quasiconformal deformations, quadratic-like maps); remarkable
functional equations (linearization, Siegel disks, Böttcher coordinates,
Écalle--Voronin cylinders); the parameter plane, combinatorics of external rays,
Thurston theory; and little Mandelbrot copies, the straightening theorem and the
Yoccoz puzzle.

\subsection*{Definitions}

\begin{definition}[Basin of infinity, filled Julia set, Mandelbrot set]
\label{def:lyubich-basic}
For $f_c(z)=z^2+c$ let $D(\infty)=\{z:\ f_c^n(z)\to\infty\}$.  Then
$K(f_c)=\widehat{\C}\setminus D(\infty)$, $J(f_c)=\partial K(f_c)$, and
\[
 \M=\{c:\ K(f_c)\text{ is connected}\}=\{c:\ f_c^n(0)\not\to\infty\}.
\]
\end{definition}

\begin{definition}[Green function, Böttcher coordinate, external rays]
\label{def:lyubich-green}
On the basin of infinity the Green function is
$G_c(z)=\lim_{n\to\infty}2^{-n}\log^+|f_c^n(z)|=\log|B_c(z)|$, where $B_c$ is the
Böttcher coordinate.  Equipotentials are the level sets of $G_c$.  External rays are
$R_c(\theta)=B_c^{-1}\{re^{2\pi i\theta}:r>1\}$.
\end{definition}

\begin{definition}[Quadratic-like maps, hybrid equivalence, modulus]
\label{def:lyubich-ql}
A quadratic-like map is a holomorphic degree-two branched covering $f\colon U\to U'$
between conformal disks with $U\Subset U'$, and $K(f)=\{z:\ f^n(z)\in U\ \forall n\}$,
$J(f)=\partial K(f)$.  Two quadratic-like maps are hybrid equivalent if they are
quasiconformally conjugate by a map $h$ with $\bar\partial h=0$ a.e.\ on $K(f)$.
For an annulus conformally equivalent to $\{1<|z|<R\}$ the modulus is
$\modulus(A)=\frac{1}{2\pi}\log R$.  The fundamental annulus of $f$ is $U'\setminus\overline U$.
\end{definition}

\begin{definition}[Renormalization; primitive and satellite]
\label{def:lyubich-renorm}
If $f^p$ restricts to a quadratic-like map on a smaller domain, the resulting
quadratic-like germ, after adjustment and rescaling, is the renormalization $Rf$.  It
is primitive when the associated little Julia sets are disjoint and satellite otherwise.
\end{definition}

\begin{definition}[Combinatorial rotation number]
\label{def:lyubich-rotation}
If a periodic orbit of external rays is cyclically permuted with $q$ rays, and the
permutation advances by $p$ positions, its combinatorial rotation number is $p/q\in\Q/\Z$.
\end{definition}

\begin{definition}[Puzzles and parapuzzles]
\label{def:lyubich-puzzle}
A puzzle is the partition of the dynamical plane cut by finitely many external rays,
equipotentials, and their pullbacks.  A parapuzzle is the corresponding partition of
the parameter plane, given by the same ray and equipotential combinatorics.
\end{definition}

\begin{remark}
The notes use the phrase a priori bounds for a uniform lower bound on the moduli of
fundamental annuli in successive renormalizations, but they give no standalone
definition or theorem with that name.  The relevant discussion is incomplete in the notes.
\end{remark}

\subsection*{Results}

\begin{theorem}[Lyubich notes, Theorem~4.13]
\label{thm:lyubich-4-13}
For a quadratic polynomial, $J$ and $K$ are either connected or Cantor sets.  The
Cantor case occurs exactly when $f^n(0)\to\infty$.  Status: proved in the notes.
\end{theorem}

\begin{theorem}[Lyubich notes, Theorem~6.10; Douady--Hubbard connectivity]
\label{thm:lyubich-6-10}
$\M$ is connected, and the parameter Böttcher map
$\Phi_\M\colon\C\setminus\M\to\C\setminus\overline{\mathbb D}$ is conformal and tangent to
the identity at infinity.  Status: proved in the notes.
\end{theorem}

\begin{theorem}[Lyubich notes, Theorem~6.12; multiplier theorem]
\label{thm:lyubich-6-12}
For every hyperbolic component $H\subset\operatorname{int}\M$ the multiplier map
$\lambda\colon H\to\mathbb D$ is a conformal isomorphism.  Status: proved in the notes.
\end{theorem}

\begin{theorem}[Lyubich notes, Theorem~9.2; straightening]
\label{thm:lyubich-9-2}
Every quadratic-like map is hybrid conjugate to some $f_c$; if its Julia set is
connected, $f_c$ is unique.  Status: existence and uniqueness proved in the notes.
\end{theorem}

\begin{theorem}[Lyubich notes, Theorem~7.36; main cardioid]
\label{thm:lyubich-7-36}
$\M$ is locally connected at every point of the main cardioid.  Status: proved in the notes.
\end{theorem}

\begin{theorem}[Lyubich notes, Theorems~11.5--11.6; Yoccoz]
\label{thm:lyubich-11-5}
If every periodic point of $f$ is hyperbolic and $f$ is not infinitely renormalizable,
then $J(f)$ is locally connected.  In particular this holds when both fixed points are
repelling and $f$ is not renormalizable.  Status: the notes present the puzzle proof of
Chapter~11.
\end{theorem}

\begin{theorem}[Lyubich notes, Theorem~10.15]
\label{thm:lyubich-10-15}
The interior of a little Mandelbrot copy is canonically homeomorphic to the interior of
$\M$.  Status: proved in Chapter~10.
\end{theorem}

\begin{remark}[Status of MLC in the notes]
The notes prove local connectivity of $\M$ on the main cardioid and state local
connectivity at non-renormalizable parameter points.  They do not prove global
local connectivity of $\M$, which the notes treat as open.
\end{remark}

\begin{remark}[Numbering]
The printed numbering of the notes is not fully consistent: Chapter~6 starts with
the number 28 and Chapter~10 uses numbers 42--44.  The numbers above are those
printed at the places cited, and should be checked against the printed text.
\end{remark}

\section{Exact examples from the monographs}
\label{sec:digest-examples}

The following two examples are standard in Milnor's lectures.  Both are proved
formally in the repository.

\begin{theorem}[Filled and Julia sets for $c=0$]
\label{thm:example-zero}
$K(f_0)=\{|z|\le1\}$ and $J(f_0)=\{|z|=1\}$.
\end{theorem}

\begin{proof}
The orbit of $z$ is $z^{2^n}$.  If $|z|\le1$ it stays in the closed disk, and if
$|z|>1$ then $|z|^{2^n}\ge|z|^n\to\infty$.  So $K(f_0)$ is the closed disk.  Its boundary
is the unit circle, and $J=\partial K$ by Milnor's Lemma~17.1.
\end{proof}

\begin{theorem}[Filled and Julia sets for $c=-2$]
\label{thm:example-minus-two}
$K(f_{-2})=J(f_{-2})=\{x\in\R:\ -2\le x\le2\}$.
\end{theorem}

\begin{proof}
If $x\in[-2,2]$, the real orbit stays in $[-2,2]$, since $y^2-2\in[-2,2]$ whenever
$y\in[-2,2]$.  For the converse, let $z\notin[-2,2]$.  Solve $w^2-zw+1=0$.  Its roots $w$
and $w^{-1}$ satisfy $z=w+w^{-1}$.  If one root had modulus $1$, then $z=w+\overline w$
would be real with $|z|\le2$, which is excluded.  So some root $w$ has $|w|>1$.  By
induction $f_{-2}^n(z)=w^{2^n}+w^{-2^n}$, because
$(a+a^{-1})^2-2=a^2+a^{-2}$.  Hence $|f_{-2}^n(z)|\ge|w|^{2^n}-1\to\infty$.  Thus
$K(f_{-2})=[-2,2]$.  The segment is compact with empty interior in $\C$, so it equals its
own frontier, and $J(f_{-2})=\partial K(f_{-2})=[-2,2]$.
\end{proof}

\begin{remark}
Theorem~\ref{thm:example-minus-two} places $c=-2$ at the left end of the real slice of
$\M$, which by Theorem~\ref{thm:real-slice} is $[-2,\tfrac14]$.  The parameter $c=0$ is
interior to that slice.
\end{remark}

\paragraph{Lean correspondence.}
Theorem~\ref{thm:example-zero} is
\leanref{Mlc/Quadratic/Examples.lean}{filledJulia_zero} and
\leanref{Mlc/Quadratic/Examples.lean}{julia_zero}.  Theorem~\ref{thm:example-minus-two} is
\leanref{Mlc/Quadratic/Examples.lean}{filledJulia_neg_two} and
\leanref{Mlc/Quadratic/Examples.lean}{julia_neg_two}.  The Joukowski step is
\leanref{Mlc/Quadratic/Examples.lean}{exists_joukowski_root_outside}, and the
unboundedness estimate is \leanref{Mlc/Quadratic/Examples.lean}{not_boundedOrbit_neg_two_of_root}.
